\section{Discussion}\label{sec:discussion}
\textbf{OCEAN traits can be transferred into an LLM, and the resulting adapters compose.} Across the Big Five, LoRAs trained through the constitution-guided DPO+SFT pipeline of \Cref{sec:methods} move some target traits monotonically on both the \texttt{TRAIT} benchmark and our judge panel, while leaving \texttt{MMLU} close to baseline across a broad scaling window, whilst others prove harder to move.
Inversion by negative scaling works in some traits, as do linear combinations of trait LoRAs, to produce approximately the expected trait mixtures (\Cref{fig:intro-main,fig:combination-heatmap}), and the LoRA reaches the same level of trait expression at higher coherence than activation capping.

\textbf{The trait axes carry downstream behavioural signal.} The more load-bearing finding for safety is that movement along an OCEAN axis transports predictable real-world behaviour with it.
Suppressing neuroticism flattens the per-turn ``frustration'' trajectory of \citet{soligo2026gemma}; moving along the agreeableness axis trades off sycophantic axis (\Cref{sec:downstream}).
 None of these benchmarks appear in training --- and they suggest that OCEAN is picking up genuine behavioural structure inside the model rather than just a rubric for a questionnaire.
At the same time, the agreeableness/sycophancy coupling shows that safety-relevant axes may not be orthogonal to the personality axes we think we are steering, which is both an opportunity for alignment and a caution about off-target effects.

\textbf{LLM personas are adjacent to, but not identical to, human personality structure.} Inter-trait correlations in our LoRAs partly mirror the correlations in human populations (e.g.\ conscientiousness and neuroticism covary in the same direction as in \citet{digman1990personality}) and partly do not.
We read this as consistent with the broader thesis of the paper: human psychometrics is a useful starting basis for studying LLM behaviour because it comes with validated instruments, but the goal of persona research on LLMs is eventually to move beyond the human frame and characterise the model's own persona space.
The unsupervised work in \Cref{sec:unsupervised-persona-exploration} is a first step in that direction; the fact that we can recover an injected ``disagreeable'' signal without telling the pipeline about OCEAN is encouraging, even if the unassisted factors are still noisy.

% \textbf{Treating personas as `weights + context' allows us to apply psychometric methods to LLM behaviour.}
% Using the framework laid out in \Cref{sec:unsupervised-persona-exploration}, we are able to discover latent behavioural factors which explain a large amount of variance in model behaviour.

\subsection{Limitations}\label{sec:limitations}

\textbf{Models selection.} The main results were produced on Llama-3.1-8B-Instruct, with a smaller set of validation runs on Gemma-3 27B. We would like to see this study replicated over different model families and sizes. Whether the same trait axes, adapter geometry, and downstream behavioural couplings persist in larger open-weight models, in different families (e.g.\ Qwen, Mistral) remains an open question.

\textbf{Evaluation and measurement.} Our evaluation stack mixes a psychometric MCQ benchmark (\texttt{TRAIT}), capability benchmarks (\texttt{MMLU}, \texttt{GSM8K}, \texttt{TruthfulQA}), and LLM-judge panel. This is broader than many prior reports, but it is still a narrow slice of the evaluation space. More fundamentally, \texttt{TRAIT} was designed for humans, and it is not obvious that a human-calibrated questionnaire measures a well-defined construct when applied to an LLM whose persona is not stable across contexts. We may also want to look at variation in personas in multilingual settings, or over long-context agentic coding settings.

\textbf{Synthetic exploration of traits.} The analysis presented in \Cref{sec:unsupervised-persona-exploration} is exploratory work using highly diverse but nonetheless synthetic rollouts. A study using real deployment trajectories across a range of tasks, and using a range of models, would enable greater confidence in the methodology and possibly more interesting latent dimensions. Validation of trait-induction was performed using the questionnaire which found the factors, rather than independent judges as discussed above. Nevertheless, we believe this area is ripe for further work.

% \textbf{Interpretation of weight-space representation of traits.} \sid{This work does not perform an in-depth analysis of the relationship between the weights we train for different traits (e.g. same-trait amplifier vs. suppressor), nor into the representation of traits in the weights themselves (where weights exist, across different models). Such a study could improve our understanding of how model personas are shaped mechanistically, and how this can affect behaviour.}

\section{Related work}\label{sec:related-work}

\begin{itemize}
  \item Assistant axis and anthropics persona vectors for thinking about personas at all, and that work has been looking at these in activation space
  \item Open Character Training showing that we can fine tune LoRas to amplify traits
  \item Papers around OCEAN
  \item Anthropic Persona Selection Model
  \item Activation-Space Personality Steering: Hybrid Layer Selection for Stable Trait Control in LLMs
  \item Editing Models With Task Arithmetic
  \item Chat Vector: A Simple Approach to Equip LLMs with Instruction Following and Model Alignment in New Languages
  \item Rediscovering the Latent Dimensions of Personality with Large Language Models as Trait Descriptors
  \item …
\end{itemize}



% \section{Further Work}\label{sec:further-work}

The study of LLM personas is a burgeoning field with much work to be done, and we would love to see various improvements and extensions to this work. This section will not be exhaustive, but hopefully point towards low-hanging fruit that might help further validate and improve the assumptions and methodology.
\textbf{Improving the methodology for the supervised case}. Creating better LoRAs (i.e. ones with either more coverage of LLM behaviour, or with more independence allowing for cleaner combinations) would allow for more robust study of the relationship between trait space and personas. In addition, there would be benefit to further investigation of different ways to combine, and especially to invert, the trait LoRAs.
\textbf{Interpretability of the traits in weight space.} We would like to see exploration of where within a model different traits are shaped most, which could be done by further study into the weights themselves. Training LoRAs for the same trait across different models, or similar traits on the same model, could shed light into whether there is a consistent relationship between traits and weights, and whether continuity in one translates to continuity in the other.
\textbf{Further study into the relationship between trait space and personas.} Our work studying the traits we have has been limited to performing statistical analysis and visualising the flattened LoRA vectors. There could be a better (i.e. more predictive or more explainable) way of embedding this information, and we would be excited by work which seeks to discover this.
\textbf{More unsupervised exploration of personas.} From a safety perspective, understanding the non-human traits of LLMs is useful. Testing alternative methods to either naturally elicit diverse personas which exist within LLMs, or to cluster and label them, would allow better study of such behaviour. We would also be excited to see other, entirely different methods for the unsupervised exploration of personas.



\subsection{Conclusion}